\documentclass[12pt]{article}
\usepackage[margin=0.72in]{geometry}
\usepackage{lmodern}
\usepackage{microtype}
\usepackage{multicol}
\usepackage{fancyhdr}
\usepackage{titlesec}
\usepackage{enumitem}
\usepackage{xcolor}
\usepackage[hidelinks]{hyperref}

\setlength{\columnsep}{0.30in}
\setlength{\parindent}{1.05em}
\setlength{\parskip}{0.24em}
\setlength{\headheight}{15pt}
\titleformat{\section}{\large\bfseries\scshape}{}{0pt}{}
\titleformat{\subsection}{\normalsize\bfseries}{}{0pt}{}
\pagestyle{fancy}
\fancyhf{}
\fancyhead[L]{Augustine on Free Choice}
\fancyhead[R]{Reading Handout}
\fancyfoot[C]{\thepage}
\renewcommand{\headrulewidth}{0.4pt}

\newcommand{\focusbox}[1]{%
\vspace{0.4em}\noindent\colorbox{gray!12}{\begin{minipage}{0.96\linewidth}#1\end{minipage}}\vspace{0.5em}}

\begin{document}
\begin{center}
{\LARGE\bfseries Are We Free, or Are We Ruled by Fate?}\par\vspace{0.25em}
{\large\scshape Reading 3: Augustine on Free Choice and Evil}\par\vspace{0.45em}
\rule{0.82\textwidth}{0.4pt}\par\vspace{0.7em}
\end{center}

\section*{Vocabulary Preview}
\begin{multicols}{2}
\begin{itemize}[leftmargin=*, itemsep=0.18em]
\item \textbf{Free choice}: The power of the will to turn toward or away from what is good.
\item \textbf{Will}: The inner power by which a person chooses, loves, seeks, refuses, or consents.
\item \textbf{Moral evil}: Evil that is done by a person, such as sin, injustice, or disordered action.
\item \textbf{Suffering evil}: Evil that is undergone, such as punishment, pain, or loss.
\item \textbf{Sin}: A voluntary turning away from what is truly and permanently good toward lower, changeable goods.
\item \textbf{Disordered desire}: Wanting a lower good in the wrong way, or loving it more than a higher good.
\item \textbf{Eternal law}: The unchanging order of justice by which higher things should rule lower things.
\item \textbf{Temporal law}: Human law that governs earthly society and may change with circumstances.
\item \textbf{Reason}: The part of the mind that can understand truth and should rule appetite and impulse.
\item \textbf{Passion}: A strong movement of desire, fear, anger, pleasure, or grief.
\item \textbf{Mutable}: Changeable; able to be gained, lost, altered, or destroyed.
\item \textbf{Providence}: God's wise rule over creation.
\end{itemize}
\end{multicols}

\section*{Introduction}
Augustine of Hippo lived from A.D. 354 to 430. Before becoming a Christian bishop, he struggled intensely with the problem of evil. If God is good and made all things, where does evil come from? Augustine had once been attracted to Manichaeism, a religion that explained evil as if it were an opposing power against God. In \textit{On Free Choice of the Will}, written as a dialogue with Evodius, Augustine rejects that answer.

Augustine asks whether God causes evil, whether human beings are responsible for sin, and what makes a will bad. This reading gives students a Christian and late ancient answer to the unit question. Aristotle began with praise, blame, compulsion, and ignorance. Cicero asked whether causes become fate. Augustine presses deeper: if the will turns away from the highest good, can it blame God, fate, law, desire, or anything outside itself?

\focusbox{\textbf{Source note:} Adapted and modernized from Augustine, \textit{On Free Choice of the Will} (\textit{De libero arbitrio}), Book I, selected sections. This classroom version follows Augustine's sequence of argument from the Latin text: evil done and evil suffered, sin as disordered desire, eternal and temporal law, reason ruling desire, good will, and evil as a voluntary turning from eternal goods to changeable goods.}

\begin{multicols}{2}
\section*{Reading}

\subsection*{1. Is God the Cause of Evil?}
Evodius begins with a hard question: Is God the cause of evil?

Augustine answers by making a distinction. We use the word \textit{evil} in two different ways. Sometimes we mean the evil a person \textit{does}: lying, betraying, murdering, worshiping falsely, or choosing what is unjust. Sometimes we mean the evil a person \textit{suffers}: pain, loss, punishment, fear, or misery.

God, Augustine says, is not the cause of evil action. God is good, and goodness does not do evil. But God is just, and justice may punish wrongdoing. Punishment is painful to the person who suffers it, so it can be called an evil suffered. Yet if punishment is just, its source is not wickedness but justice.

So Augustine's first answer is this: God does not cause moral evil. Each wrongdoer is the cause of his own wrongdoing. If evil actions were not done by the will, they could not be punished justly.

\subsection*{2. Evil Is Not Truly Taught}
Evodius then wonders whether people do evil because they have learned to do evil. Augustine presses him: is teaching itself good or bad?

Evodius admits that teaching is good, because teaching leads to knowledge and understanding. But understanding is not evil. Therefore, Augustine says, evil is not truly learned as a skill to be practiced. If we learn about evil, we learn that it should be avoided.

A real teacher teaches what is good. To do wrong is not to obey teaching, but to turn away from it. Augustine is preparing an important point: sin does not come from knowledge as knowledge. It comes from a will that refuses the order it knows or should know.

\subsection*{3. Look Beneath the Outward Act}
Augustine next asks what makes an action wrong. Evodius names obvious sins: adultery, murder, sacrilege, and similar acts. But Augustine wants to know why they are wrong.

Evodius first appeals to human law. Augustine answers that law forbids adultery because it is wrong; adultery is not wrong merely because law forbids it. Evodius then says that adultery is wrong because no one would want it done to himself. Augustine shows that this rule is not enough. A corrupt person might willingly trade permission for such wrongdoing and still be guilty.

The evil, Augustine says, is not found only in the outward act that can be seen. The deeper evil is disordered desire. A person who desires adultery and would commit it if he had the chance is already guilty in will, even if no outward act occurs.

This helps Augustine define sin more deeply. Sin is not merely a bad movement of the body. Sin begins when the will loves wrongly.

\subsection*{4. Not Every Desire Is Evil}
At this point Augustine makes the problem harder. If sin comes from desire, is every desire evil?

Consider a man who kills because he is afraid. He may not desire money or pleasure. He desires to live without fear. But living without fear is a good thing. Therefore, not every desire is blameworthy.

The problem is not desire itself, but disordered desire. The good person wants freedom from fear by turning away from goods that can be lost against his will. The wicked person wants freedom from fear so that he may safely cling to lower things: power, pleasure, revenge, or possessions.

Augustine's point is severe: the same outward goal may hide different loves. One person seeks peace by ordering his soul. Another seeks peace by removing anyone who blocks his appetite.

\subsection*{5. Temporal Law and Eternal Law}
Augustine then distinguishes two kinds of law.

Temporal law is the law of earthly society. It protects public order, property, safety, families, and civic peace. Because human societies change, temporal law may also change. A disciplined people may be trusted with one form of rule; a corrupt people may need another. A human law can be just in one circumstance and later need to be changed.

But temporal law depends on a higher law. Augustine calls this eternal law. Eternal law does not change. It is the order of justice itself: higher things should rule lower things, reason should rule appetite, good should be loved more than what is merely useful or pleasant, and all things should be rightly ordered.

Human law can punish theft, violence, and public crimes. But it cannot cure every disordered love. A person may obey the law outwardly and still be ruled inwardly by greed, pride, lust, envy, or fear. Therefore the deepest judgment of the soul belongs not merely to human law, but to the eternal order of justice.

\subsection*{6. What Should Rule in a Human Being?}
Augustine asks what makes human beings higher than animals. It is not strength, speed, senses, appetite, or bodily life. Many animals surpass us in those things. Human beings are higher because they have reason and understanding.

Therefore a well-ordered person is one in whom reason rules the lower movements of the soul. Appetite, anger, fear, pleasure, ambition, and grief are not meant to rule. They belong under the guidance of reason.

When reason rules desire, the human person is rightly ordered. When desire rules reason, the person is disordered. Augustine calls this slavery, even if the person seems outwardly free. A man may have no human master and still be enslaved by his own passions.

Could desire force a good mind to become its slave? Augustine says no. Lower things cannot overpower a mind that is rightly ordered by virtue. Equal or higher things would not unjustly overthrow it. Therefore nothing makes the mind a companion of disordered desire except its own will and free choice.

\subsection*{7. The Good Will}
Evodius worries that if foolish people suffer misery, perhaps they are punished for a wisdom they never had. Augustine turns the discussion to something every person can recognize: the will.

Do we have a will? Augustine asks. Evodius hesitates, but Augustine points out that even asking a question assumes a will to know. Friendship, learning, wisdom, and happiness all involve willing.

Then Augustine asks whether we can have a good will. A good will is the will by which we desire to live rightly and honorably and to reach wisdom. This good will is greater than riches, honors, bodily pleasures, or public glory, because those things can be lost against our will. But the will to love what is right is within us.

If a person truly loves a good will more than changeable goods, the virtues begin to appear. Prudence sees what should be sought and avoided. Fortitude can bear the loss of things not in our control. Temperance resists shameful desire. Justice gives each person what is due. A good will is not a small possession. Augustine sees it as the beginning of a rightly ordered life.

\subsection*{8. Why Do All Want Happiness but Not All Have It?}
Everyone wants to be happy. Augustine does not deny this. But not everyone wants to live rightly. This explains why many people want happiness and yet remain miserable.

A wicked person wants the reward without the order that deserves it. He wants happiness while clinging to the very loves that make the soul restless and enslaved. Eternal law joins happiness to right willing and misery to disordered willing. The miserable do not choose misery as their goal, but they choose a kind of will that necessarily brings misery with it.

So Augustine's answer is not that human beings are free because nothing influences them. His answer is that the will has a real direction. It either turns toward the order of truth and goodness, or it turns toward changeable things as if they were ultimate.

\subsection*{9. Augustine's Definition of Sin}
At the end of Book I, Augustine returns to the original question: what is it to do evil?

Doing evil means neglecting eternal goods and chasing temporal goods as if they were the highest things. Eternal goods are loved without fear of losing them against one's will: truth, wisdom, justice, rightly ordered love, and God. Temporal goods are not evil in themselves. Bodies, families, honors, money, food, beauty, and civic life all have their proper place. The problem is not that these things exist. The problem is that a disordered soul clings to them as if they could make it finally secure.

Evodius accepts the conclusion: all sins are contained in this pattern. The soul turns away from divine and lasting things and turns toward changeable and uncertain things. The things themselves may be good in their order, but the soul becomes perverse when it submits to what it ought to rule.

Augustine therefore concludes that evil comes from free choice of the will. But this raises the next problem: if free choice gives us the power to sin, should God have given it to us at all? Augustine says that question must be taken up next. For now, Book I has reached its answer: God is not the cause of sin. The will becomes guilty when it freely turns from higher good to lower good.

\section*{Reading Focus}
\begin{enumerate}[leftmargin=*]
\item What is Augustine's distinction between evil that is \textit{done} and evil that is \textit{suffered}?
\item Why does Augustine say the evil in sin is not found only in the outward act?
\item What is the difference between temporal law and eternal law?
\item Augustine says reason should rule desire. How does this connect to Aristotle's claim that habit and character shape action?
\item Why can everyone want happiness while not everyone receives it, according to Augustine?
\end{enumerate}

\section*{Window Note and Summary}
Create a window note showing Augustine's two directions of the will: turning toward eternal goods and turning toward changeable goods. Include at least one drawing or symbol for each direction. At the bottom, write a 3--5 sentence summary explaining why Augustine thinks God is not the cause of moral evil.

\section*{Connection to the Unit Question}
Write one claim answering this question: \textbf{If our desires pull us toward lower goods, are we still responsible for what we choose?} Use Augustine to support your answer, and make one comparison to either Aristotle or Cicero.
\end{multicols}
\end{document}
